\section*{Abstract}
		The T0-Time-Mass-Duality theory derives fundamental constants and masses parameter-free from the universal geometric parameter $\xi = 4/30000$. This complementary document examines a validation of the fractal dimension $D_f = 3 - \xi \approx 2.99987$ through backward derivation from the experimental mass ratio $r = m_{\mu} / m_e \approx 206.768$ (CODATA 2025). The direct geometric calculation chain used here does not hold in its present form: the electron formula gives $0.0854$~GeV instead of $m_e$, the muon chain gives $m_e$ instead of $m_\mu$, and the Nambu inversion is not consistent with the torsion mass $m_T$; a validation of $D_f$ does not follow from it. The systematic mass calculation is presented in \emph{006\_T0\_Teilchenmassen\_En.pdf}.

	
	
	
	\section{Introduction}
	\label{sec:introduction}
	
	\begin{important}{Document Complementarity}{}
		This document focuses on \textbf{examining a validation of the fractal dimension} $D_f$ from experimental lepton masses. It complements the main document \emph{006\_T0\_Teilchenmassen\_En.pdf}, which presents the complete systematic mass calculation for all fermions.
	\end{important}
	
	Particle physics faces the fundamental problem of arbitrary mass parameters in the Standard Model. The T0-Time-Mass-Duality theory revolutionizes this approach through a completely parameter-free description.
	
	\section{Parameters and Basic Formulas}
	\label{sec:parameters}
	
	The theory is based on time-energy duality and fractal spacetime structure.
	
	\subsection{Exact Geometric Parameters}
	\label{subsec:exact_parameters}
	
	\begin{align}
		\xi &= \frac{4}{30000} = \frac{1}{7500} \approx 1.333 \times 10^{-4}, \label{eq:xi} \\
		D_f &= 3 - \xi \approx 2.99986667, \label{eq:Df} \\
		\alpha &= \frac{1 - \xi}{137} \approx 7.298 \times 10^{-3}, \label{eq:alpha} \\
		K_{\text{frac}} &= 1 - 100 \xi \approx 0.9867, \label{eq:K} \\
		g_{T0}^2 &= \alpha K_{\text{frac}}, \label{eq:gT0} \\
		E_0 &= \frac{1}{\xi} \approx \SI{7500}{\giga\electronvolt}, \label{eq:E0} \\
		p &= -\frac{2}{3}. \label{eq:p}
	\end{align}

	$1/\xi = 7500$ is a pure number; assigning the unit GeV in Eq.~\eqref{eq:E0} is a setting, not a derivation (other documents read the same number as eV). In T0 units with $\alpha = 1$ one has $\xi E_0^2 = 1$, i.e.\ $E_0^2 = 1/\xi = 7500$; the number belongs to $E_0^2$, not to $E_0$ (Doc.~386). The corpus value is $E_0 = \sqrt{m_e m_\mu/K_{\text{frac}}} = 7.398$~MeV; $7500$~GeV exceeds it by a factor of $1.01\times10^{6}$. The following formulas use $E_0\,\xi = 1$~GeV and depend on this setting.
	
	\begin{result}{Fine Structure Constant Precision}{}
		The deviation of $\alpha$ from CODATA is only $\approx 0.013\%$ -- strong evidence for the fractal correction.
	\end{result}
	
	\section{Geometric Mass Derivation - Direct Method}
	\label{sec:geometric_derivation}
	
	T0-theory offers several mathematically equivalent methods for mass calculation. In this document the \textbf{direct geometric method} is examined as to whether it validates the fractal dimension.
	
	\subsection{Electron Mass $m_e$ - Direct Geometric Method}
	\label{subsec:electron_mass}
	
	In the direct geometric method:
	\begin{align}
		m_e &= E_0 \cdot \xi \cdot \sqrt{\alpha} \cdot \frac{\Gamma(D_f)}{\Gamma(3)}. \label{eq:me_direct}
	\end{align}

	With the parameters of Section~\ref{subsec:exact_parameters} ($E_0\,\xi = 1$~GeV, $\sqrt{\alpha} = 0.08543$, $\Gamma(2.99987)/\Gamma(3) = 0.99994$) the formula gives $m_e = 0.0854$~GeV -- larger than the target value $5.10\times10^{-4}$~GeV by a factor of $167$. The formula does not yield $m_e$ in this form. The deviation of $-0.20\%$ from CODATA ($\SI{0.000511}{\giga\electronvolt}$) and the following table refer to the target value, not to the result of the formula.
	
	\subsection{Consistency Check with Main Document}
	\label{subsec:consistency_check}
	
	\begin{table}[H]
		\centering
		\resizebox{\textwidth}{!}{
\begin{tabular}{lccc}
			\toprule
			\textbf{Method} & \textbf{$m_e$ [GeV]} & \textbf{Accuracy} & \textbf{Source} \\
			\midrule
			Direct geometric (target value) & $5.10\times10^{-4}$ & $99.8\%$ & This document \\
			Extended Yukawa & $5.11\times10^{-4}$ & $99.9\%$ & 006 \\
			Experiment (CODATA) & $5.11\times10^{-4}$ & $100\%$ & Reference \\
			\bottomrule
		\end{tabular}
}
		\caption{Consistency of mass calculation methods in T0-theory}
		\label{tab:method_consistency}
	\end{table}
	
	\begin{result}{Method Comparison}{}
		The agreement within $0.2\%$ holds for the target value of the direct method, not for the result of Eq.~\eqref{eq:me_direct}; an equivalence of the methods is therefore not shown. The Yukawa method bridges to the Standard Model.
	\end{result}
	
	\subsection{Effective Torsion Mass $m_T$}
	\label{subsec:torsion_mass}
	
	\begin{align}
		R_f &= \frac{\Gamma(D_f)}{\Gamma(3)} \sqrt{\frac{E_0}{m_e}}, \label{eq:Rf} \\
		m_T &= \frac{m_e}{\xi} \sin(\pi \xi) \, \pi^2 \sqrt{\frac{\alpha}{K_{\text{frac}}}} \, R_f \approx \SI{5.220}{\giga\electronvolt}. \label{eq:mT}
	\end{align}
	
	\subsection{Muon Mass $m_{\mu}$}
	\label{subsec:muon_mass}
	
	From RG-duality and loop integral $I$:
	\begin{align}
		I &= \int_0^1 \frac{m_e^2 x (1-x)^2}{m_e^2 x^2 + m_T^2 (1-x)}  dx \approx 1.60 \times 10^{-9}, \label{eq:I} \\
		r &\approx \sqrt{6 I} \approx 9.8 \times 10^{-5}, \label{eq:r} \\
		m_T \cdot r &\approx \SI{5.11e-4}{\giga\electronvolt}. \label{eq:mmu}
	\end{align}

	The integral is evaluated numerically with $m_e = 5.11\times10^{-4}$~GeV and $m_T = 5.220$~GeV. The product $m_T\,r$ thus gives $m_e$, not $m_\mu$ ($\SI{0.105658}{\giga\electronvolt}$); the chain does not yield $m_\mu$ in this form.
	
	\begin{important}{Mass Ratio}{}
		Since the chain does not yield $m_\mu$, no mass ratio $r = m_{\mu} / m_e$ can be calculated from it; an independent validation of the geometric foundation does not result here.
	\end{important}
	
	\section{Backward Validation: $D_f$ from $r$ and Nambu Formula}
	\label{sec:backward_validation}
	
	The classical Nambu formula $r \approx (3/2)/\alpha$ (dev. $-0.58\%$) is refined by the $\xi$-correction.
	
	\subsection{Nambu Inversion}
	\label{subsec:nambu_inversion}
	
	\begin{align}
		m_T^{\text{target}} &= \frac{m_{\mu}}{\sqrt{\alpha} \cdot (3/2) \cdot (1 - \xi)} = \frac{0.10566}{0.08543\times1.5\times0.99987} \approx \SI{0.825}{\giga\electronvolt}. \label{eq:mTtarget}
	\end{align}

	This value is not consistent with Equation~\eqref{eq:mT} ($m_T = 5.220$~GeV); the inversion does not hold in this form.
	
	\subsection{Optimization for $D_f$}
	\label{subsec:optimization_df}
	
	Define $m_T(D_f)$ according to Equation~\ref{eq:mT} and solve:
	\begin{align}
		D_f = \arg\min \left| m_T(D_f) - m_T^{\text{target}} \right|. \label{eq:optDf}
	\end{align}
	
	\begin{keyresult}{Result of the Backward Validation}{}
		Since $m_T^{\text{target}} = 0.825$~GeV and $m_T = 5.220$~GeV do not agree, the optimization does not hold on this basis; a result $D_f \approx 2.99986667$ is not supported by this calculation. That the experimental mass ratio compels the fractal geometry does not follow from this calculation; the foundations of \emph{006\_T0\_Teilchenmassen\_En.pdf} are not independently validated here.
	\end{keyresult}
	
	\section{Application: Anomalous Magnetic Moment $a_{\mu}^{\text{T0}}$}
	\label{sec:application_g2}
	
	With the fractal dimension $D_f = 3 - \xi$ and geometric masses:
	\begin{align}
		F_2^{\text{T0}}(0) &= \frac{g_{T0}^2}{8 \pi^2} I_{\mu} K_{\text{frac}}, \label{eq:F2} \\
		\text{term} &= \left( \frac{\xi E_0}{m_T} \right)^p = m_T^{2/3}, \label{eq:term} \\
		F_{\text{dual}} &= \frac{1}{1 + \text{term}} \approx 0.249, \label{eq:Fdual} \\
		a_{\mu}^{\text{T0}} &= F_2^{\text{T0}}(0) \cdot F_{\text{dual}} \approx 1.53 \times 10^{-9} = 153 \times 10^{-11}. \label{eq:amu}
	\end{align}
	
	\begin{result}{Experimental Validation}{}
		Status 2025 (Fermilab final result, 2025 White Paper of the Muon~$g-2$ Theory Initiative): $\Delta a_\mu = (38 \pm 63) \times 10^{-11}$. The T0 value $153 \times 10^{-11}$ lies $(153 - 38)/63 \approx 1.8\sigma$ above it.
	\end{result}
	
	\section{Python Implementation and Reproducibility}
	\label{sec:python_implementation}
	
	\begin{important}{Full Transparency}{}
		For reproduction of all numerical calculations see the external script \texttt{t0\_df\_from\_masses\_geometry.py} in the repository folder.
	\end{important}
	
	\section{References}
	\label{sec:references}
	
	\begin{itemize}
		\item Pascher, J. (2025). \emph{T0-Model: Complete Parameter-Free Particle Mass Calculation} (006\_T0\_Teilchenmassen\_En.pdf). Available at: 
		
		\item Pascher, J. (2025). \emph{T0-Time-Mass-Duality Repository}, GitHub v1.6. Available at: 
		
		\item CODATA (2025). \emph{Fundamental Physical Constants}, NIST.
	\end{itemize}
